\chapter{Colours}
\label{ch:colours}

Four colours --- \emph{schwarz}, \emph{wei\ss{}}, \emph{rot}, \emph{blau} ---
and a small surprise inside them. German's colour words are all home-grown
Germanic, with no Latin borrowed in; and two of them travelled the other way,
out of Germanic and into the Romance languages.
\emph{Practice lessons: \texttt{lessons/GE-C13-*.md}.}

\section{\emph{schwarz}, \emph{wei\ss{}} --- black and white, and the word German lent away}
\label{lesson:schwarz-weiss}

German's colour words are \textbf{home-grown} --- no Latin here. But one of them
has a sibling that emigrated: a German word for ``shining'' became the word for
\textbf{white} in French, Italian and Portuguese.

\begin{sounds}
\emph{schwarz} $=$ \emph{SHVARTS}: the \emph{schw-} cluster is ``sh'' $+$ ``v,''
and the final \emph{z} is always ``ts.''\\[3pt]
\emph{wei\ss{}} $=$ \emph{vice}: German \emph{w} is English ``v,'' \emph{ei} is
the vowel of English \emph{eye}, and the \ss{} (\emph{Eszett}) is a sharp ``ss.''
\end{sounds}

\begin{cousinweb}
\textbf{schwarz} (``black'') $\leftarrow$ Proto-Germanic \emph{*swartaz}. Its
English cousin survives in \textbf{swarthy} (``dark-complexioned''). Further
back it is related to Latin \emph{sord\=es} (``dirt'') --- the root of English
\textbf{sordid}. Black and grubby, from one idea.
\end{cousinweb}

\begin{cousinweb}
\textbf{wei\ss{}} (``white'') $\leftarrow$ Proto-Germanic \emph{*hw\=\i taz} ---
\emph{exactly} English \textbf{white}. The same word in different spelling
clothes.
\end{cousinweb}

\begin{sounds}
The \ss{} (\emph{Eszett}) marks a sharp \emph{s} after a long vowel. So
\emph{wei\ss{}} rhymes with English \emph{vice}, not with the second half of
\emph{vice-versa}. In Switzerland you will see it written \textbf{weiss} --- the
letter is optional there.
\end{sounds}

\subsection*{The word that emigrated}

German also has \textbf{blank} --- ``shiny, polished, bare.'' That word, in its
Frankish form \emph{blank} (``shining''), was \textbf{borrowed into Romance},
where it pushed Latin's own \emph{albus} out and became:

\begin{center}
\begin{tabular}{@{}ll@{}}
\toprule
Language & ``white'' \\
\midrule
French & \textbf{blanc} \\
Italian & \textbf{bianco} \\
Portuguese & \textbf{branco} \\
German & \emph{blank} --- still just ``\textbf{shiny}'' \\
\bottomrule
\end{tabular}
\end{center}

\begin{culture}
This is the \textbf{reverse} of the usual flow. Latin normally lends \emph{to}
Germanic (\emph{Fenster} $\leftarrow$ \emph{fenestra}, \emph{Wein}
$\leftarrow$ \emph{v\=\i num}). Here Germanic lent \emph{to} Latin's daughters
--- and German kept the original meaning while they turned it into a colour.
\end{culture}

Practise them on something you can already name: \emph{der Wein ist wei\ss{}}
(``the wine is white''), and \emph{wei\ss{}} $\cdot$ \emph{white} said side by
side, so you can hear that it is one word.

\section{\emph{rot}, \emph{blau} --- an ancient cousin and a second emigrant}
\label{lesson:rot-blau}

Two more native German colours. One is related to the French word for red ---
genuinely, not by borrowing. The other is a \textbf{second} German word the
Romance languages took.

\begin{sounds}
\emph{rot} $=$ \emph{roht}, with a long ``o'' as in English \emph{road}.\\[3pt]
\emph{blau} $=$ \emph{blow} (rhyming with \emph{cow}): German \emph{au} is the
vowel of English \emph{house}.
\end{sounds}

\begin{cousinweb}
\textbf{rot} $\leftarrow$ Proto-Germanic \emph{*raudaz} $\leftarrow$
Proto-Indo-European \emph{*h\ensuremath{_1}rewd\ensuremath{^h}-}. That root is
one of the oldest colour words we can reconstruct, and it went everywhere:
\end{cousinweb}

\begin{center}
\begin{tabular}{@{}ll@{}}
\toprule
Language & Word \\
\midrule
German & \textbf{rot} \\
English & \textbf{red}, \emph{rust}, \emph{ruddy} \\
Latin & \emph{ruber}, \emph{rubeus} $\to$ French \textbf{rouge}, \emph{ruby},
\emph{rubric} \\
\bottomrule
\end{tabular}
\end{center}

So \emph{rot} and \emph{rouge} are \textbf{cousins}, not borrowings --- they
split apart thousands of years before either language existed. Red is the colour
humans name first, after black and white.

\begin{cousinweb}
\textbf{blau} $\leftarrow$ Proto-Germanic \emph{*bl\=ewaz}. Like \emph{blank} in
the last lesson, this word was \textbf{borrowed into Romance}.
\end{cousinweb}

\begin{center}
\begin{tabular}{@{}ll@{}}
\toprule
Language & ``blue'' \\
\midrule
German & \textbf{blau} \\
French & \textbf{bleu} \\
Italian & \textbf{blu} \\
English & \textbf{blue} --- taken from \textbf{French}, not from German \\
\bottomrule
\end{tabular}
\end{center}

\begin{culture}
English is the odd one here: it had Germanic \emph{bl\=aw} available but
borrowed the French form instead. The word went Germanic $\to$ French $\to$
English, coming home in disguise. Portuguese went a different way entirely ---
its blue is \textbf{azul}, from \textbf{Arabic}.
\end{culture}

\subsection*{The chapter in one table}

\begin{center}
\begin{tabular}{@{}lll@{}}
\toprule
German & Source & Did Romance borrow it? \\
\midrule
\emph{schwarz} & Germanic \emph{*swartaz} & no \\
\emph{wei\ss{}} & Germanic \emph{*hw\=\i taz} & no --- but \emph{blank} went
instead \\
\emph{rot} & Germanic \emph{*raudaz} & no --- Latin had its \textbf{own} cousin \\
\emph{blau} & Germanic \emph{*bl\=ewaz} & \textbf{yes} --- \emph{bleu},
\emph{blu} \\
\bottomrule
\end{tabular}
\end{center}

Say all four together --- \emph{schwarz, wei\ss{}, rot, blau} --- and then the
cousins that prove the family: \emph{rot} $\cdot$ \emph{red} $\cdot$
\emph{rouge}. Next chapter: putting these onto the things you already name.
